\documentclass{article}
\usepackage{graphicx} %pa' las imágenes
\usepackage[utf8]{inputenc} %pa' los acentitos
\usepackage{array} %para la tablita
\usepackage[spanish,es-tabla]{babel}%para que aparezcan las cositas en español
\graphicspath{{/home/Valeria/Practica01/Imagenes}}

\title{\Huge Taller de Herramientas Computacionales}
\author{\huge Valeria Ortiz Cervantes}
\date{\LARGE 2 de marzo del 2019}

\begin{document}
\maketitle
\begin{center}
	\subsection*{\LARGE Universidad Nacional Autónoma de México.\\Facultad de Ciencias.\\}
	\includegraphics[scale=3]{/1.jpg}
\end{center}
\newpage
\begin{center}
	\title{\LARGE Práctica 1.Preguntas}
\end{center}
\section*{Historia del Paradigma Orientado a Objetos.}
En los años 60, en el Centro de Cálculo Noruego, un equipo de investigación, liderado por Krinsten Nygaard, tuvo problemas en el desarrollo de simulaciones de sistemas físicos, pues los programas eran muy complejos y las modificaciones, al ser una simulación, eran muy frecuentes. La solución que idearon fue diseñar el programa paralelamente al objeto físico, es decir, si el objeto físico tenía treinta componentes, el programa también tendría treinta módulos. Seccionando el programa de esta manera, había una total correspondencia entre el sistema físico y el sistema informático.\\Además de resolverse las dificultades anteriores, se obtuvo otro beneficio: la reusabilidad, poder guardar cosas para futuros programas. Para implementar esto crearon el lenguaje Simula 67.\\Después, en los años 70 en Estados Unidos, Alan Kay estaba creando un ordenador que pudiera ser utilizado por niños, Dynabook, pero también le resultó muy complicado porque era una programación nueva y experimental, por lo cual creó un lenguaje llamado Smalltalk.\\Smalltalk tuvo una gran difusión y cuando en los años 80 en ATT-Bell quisieron crear un sucesor al lenguaje C, incorporaron las principales ideas de Smalltalk y de Simula, creando el lenguaje C++. Puede afirmarse que se debe a este último la gran extensión de los conceptos de la orientación a objetos.
\begin{center}
	\section*{Historia de Python}
\end{center}
La historia de Python como lenguaje de programación inicia a finales de los 80s y principios de los 90s con Guido Van Rossum. En una navidad de 1989, Guido Van Rossum, quien trabajaba en el CWI, un centro de investigación en Ámsterdam, decidió empezar un proyecto como pasatiempo dándole continuidad a ABC, un lenguaje de programación que se desarrolló ahí mismo a principios de los 80s como alternativa a BASIC, fue pensado para principiantes por su facilidad de aprendizaje y uso. Su código era compacto pero legible.\\El proyecto no trascendió ya que el hardware disponible en la época hacía difícil su uso. Así que Van Rossum le dio una segunda vida creando Python. En 1991, Van Rossum publicó el código, en el que ya se tenían disponibles clases con herencias, manejo de excepciones, funciones y tipos modulares.\\ Además en este lanzamiento inicial aparecía un sistema de módulos adoptado de Modula-3; van Rossum describe el módulo como "uno de las mayores unidades de programación de Python". En el año 1994 se formó el foro de discusión principal de Python, marcando un hito en el crecimiento del grupo de usuarios de este lenguaje.\\
\begin{center}
	\section*{Paradigmas de programación.}
\end{center}
\subsection*{Paradigma declarativo.}
\begin{enumerate}
	\item El programa se describe en términos de proposiciones y afirmaciones que son declaradas para describir el problema.
	\item No se especifican los pasos para resolverlo.
	\item Todos los tipos de datos son inmutables.
	\item Todo está dado por funciones o expresiones puramente matemáticas.
\end{enumerate}
\subsection*{Paradigma imperativo.}
\begin{enumerate}
	\item El programa se describe en términos de instrucciones, condiciones y pasos.
	\item Se puede modificar el estado de un programa pues se permite la mutación de variables.
	\item Utilizan procedimientos, rutinas o funciones impuras para establecer mecanismos de control.
	\item Los mecanismos de control generan efectos secundarios y son quienes mutan el estado del programa durante la ejecución.
\end{enumerate}
\subsection*{Paradigma funcional.}
\begin{enumerate}
	\item Se basan en el uso de una o más funciones dentro de las cuales se pueden utilizar funciones creadas anteriormente.
	\item Su objetivo es dividir el programa en módulos de forma que cada uno de éstos realice una única función.
	\item Existen dos tipos de lenguajes funcionales:
	\begin{enumerate}
		\item Los puros, como HASKELL.
		\item Los híbridos, como SAP, ML, Scheme.
	\end{enumerate}
\end{enumerate}
\subsection*{Paradigma lógico.}
\begin{enumerate}
	\item La programación lógica comprende la programación declarativa y la funcional.
	\item Se basa en la lógica de primer orden.
	\item Especifica qué debe hacer el programa y no cómo hacerlo.
	\item Se emplea en aplicaciones de inteligencia artificial.
	\item  El lenguaje de programación lógica por excelencia es PROLOG.
\end{enumerate}
\end{document}